\documentclass[11pt]{article}
\usepackage[utf8]{inputenc}
\usepackage[T1]{fontenc}
\usepackage{natbib}
\usepackage{geometry}
\geometry{margin=1in}
\usepackage{parskip}
\usepackage{changepage}
\setlength{\parindent}{0pt}
\setlength{\parskip}{6pt}
\title{Dopaminergic Neuron Susceptibility in Long COVID: A Hypothesis on Genetic Surface Variability and SARS-CoV-2 Infection}
\author{Timothy J. VanZeben \\ Independent Researcher \\ tvanzeben@unm.edu}
\date{March 5, 2025}

\begin{document}
\maketitle
\begin{center}
\textbf{Abstract}
\end{center}

\begin{adjustwidth}{1in}{1in}
\small\itshape
Long COVID presents a diverse array of neurological symptoms \citep{taquet2021neuropsychiatric, reiken2022alzheimersignaling}, including cognitive impairment, fatigue, anhedonia, and dysautonomia. While neuroinflammation has been proposed as a key mechanism \citep{crunfli2023astrocytes, yang2021choroidplexus}, this hypothesis suggests an alternative explanation: selective infection of dopaminergic neurons due to individual genetic variability in surface receptor expression. Specifically, variants in ACE2, TMPRSS2, DRD2, and DAT1 may influence the likelihood \citep{novelli2020ace2variants, asselta2020ace2tmprss2} of SARS-CoV-2 entry into dopamine-producing cells. This could result in dopamine depletion, leading to movement disorders, motivational deficits, and executive dysfunction—symptoms that strongly correlate with Long COVID manifestations \citep{davis2021longcovidcohort, graham2021longcovidneuro}. This paper presents supporting evidence, proposes testable predictions, and discusses implications for targeted therapeutic strategies.
\end{adjustwidth}
\begin{center}
\section*{Introduction}
\end{center}

\subsection*{Background}
Long COVID affects millions of individuals worldwide, with persistent symptoms lasting weeks to months after initial SARS-CoV-2 infection. The syndrome encompasses a wide range of manifestations including severe fatigue, cognitive dysfunction ("brain fog"), autonomic dysregulation, and mood disturbances \citep{davis2021longcovidcohort}. Despite growing recognition of Long COVID as a significant public health concern, its pathophysiological mechanisms remain poorly understood.

\subsection*{Current Theories and Their Limitations}
Existing theories attribute Long COVID symptoms primarily to persistent inflammation, immune dysregulation, and microvasculature damage \citep{yang2021choroidplexus, crunfli2023astrocytes}. While these mechanisms likely contribute to symptomatology, they do not fully account for the striking resemblance between Long COVID symptoms and dopamine-related disorders such as Parkinson's disease and ADHD \citep{devine2022parkinsonscovid, gizer2009adhdgenes}. Furthermore, the high variability in symptom presentation and severity suggests additional factors may be involved in determining individual susceptibility.

\subsection*{Rationale for a Dopaminergic Hypothesis}
The neurological and cognitive symptoms of Long COVID—including fatigue, anhedonia, executive dysfunction, and motor impairments—closely parallel conditions associated with dopamine depletion \citep{brundin2020parkinsonmanagement, cappelletti2023dopamineblock}. This overlap suggests potential involvement of dopaminergic circuits in Long COVID pathophysiology. Additionally, emerging evidence indicates that SARS-CoV-2 can directly affect dopamine signaling \citep{lee2024dopamineinfection}, providing a plausible mechanism for these persistent symptoms.

\begin{center}
\section*{The Dopaminergic Susceptibility Hypothesis}
\end{center}

\subsection*{Central Proposition}
This paper proposes that SARS-CoV-2 selectively infects dopaminergic neurons in genetically susceptible individuals, leading to dopamine depletion and subsequent neurological dysfunction. The central mechanisms of this hypothesis are as follows: SARS-CoV-2 enters cells primarily through the ACE2 receptor, with TMPRSS2 facilitating viral priming. Genetic variants in ACE2, TMPRSS2, DRD2, and DAT1 may influence the susceptibility of dopaminergic neurons to viral infection. Once infected, dopaminergic neurons may experience functional impairment or undergo cell death, leading to significant dopamine depletion. This loss of dopamine signaling within specific neural circuits is proposed to underlie the diverse neurological and systemic symptoms observed in Long COVID.

\subsection*{Role of Genetic Surface Variability}
Key to this hypothesis is the substantial interindividual genetic variability in expression and function of surface proteins relevant to both viral entry and dopamine signaling. These proteins include ACE2 (Angiotensin-Converting Enzyme 2), the primary entry receptor for SARS-CoV-2 \citep{cao2020ace2genetics, hou2021ace2variants}; TMPRSS2 (Transmembrane Serine Protease 2), which facilitates viral entry through spike protein priming \citep{asselta2020ace2tmprss2}; DRD2 (Dopamine Receptor D2), a key regulator of dopamine signaling that may influence viral tropism \citep{eisenstein2016drd2}; and DAT1/SLC6A3 (Dopamine Transporter), which regulates dopamine reuptake and may affect neuronal susceptibility. These proteins are expressed in dopaminergic regions such as the substantia nigra and ventral tegmental area \citep{song2021neuroinvasion, douaud2022brain}, but their density varies significantly based on individual genetic variants. If dopamine-producing neurons in some individuals express higher ACE2 and TMPRSS2 levels, they may become prime targets for viral infection.

\begin{center}
\section*{Anatomical Distribution and Symptom Correlation}
\end{center}

\subsection*{Dopaminergic Pathway Vulnerability}
If the hypothesis is correct, symptom presentation should vary depending on which dopaminergic pathways are primarily affected. The nigrostriatal pathway (Substantia nigra $\rightarrow$ Striatum) dysfunction would manifest as muscle fatigue, motor dysfunction, and post-exertional malaise (PEM). Disruption of the mesolimbic pathway (VTA $\rightarrow$ Nucleus accumbens) would result in anhedonia, motivation loss, and post-viral depression. The mesocortical pathway (VTA $\rightarrow$ Prefrontal cortex) damage would lead to brain fog, ADHD-like symptoms, and executive dysfunction. Finally, tuberoinfundibular pathway (Hypothalamus $\rightarrow$ Pituitary gland) impairment would produce dysautonomia, sleep disturbances, and metabolic dysregulation.

\subsection*{Correlation with Observed Long COVID Symptomatology}
The symptom clusters predicted by this model align closely with reported Long COVID manifestations \citep{graham2021longcovidneuro, davis2021longcovidcohort}. For example, post-exertional malaise—a cardinal feature of Long COVID—could result from nigrostriatal disruption impairing motor function and energy regulation. Similarly, the cognitive impairment ("brain fog") frequently reported could stem from mesocortical pathway dysfunction affecting prefrontal cortex function.

\begin{center}
\section*{Supporting Evidence}
\end{center}

A growing body of evidence supports the hypothesis that SARS-CoV-2 selectively disrupts dopaminergic pathways in genetically susceptible individuals. Neuropathological studies have detected viral RNA in the substantia nigra, a region central to dopamine synthesis and motor regulation. This provides compelling evidence of direct viral neurotropism affecting dopaminergic nuclei.

Clinical presentations in long COVID patients often resemble syndromes of dopamine depletion, including psychomotor slowing, fatigue, executive dysfunction, and mood dysregulation. Neuroimaging studies corroborate these observations, showing reduced dopamine transporter (DAT) availability in post-COVID individuals. These functional deficits strongly correlate with the severity of persistent symptoms, suggesting that dopaminergic dysfunction is a central mechanism in long COVID neurocognitive sequelae.

Experimental studies confirm that SARS-CoV-2 can infect neuronal cells directly, triggering intracellular cascades that interfere with dopamine metabolism, impair mitochondrial function, and destabilize neurotransmitter regulation. These disruptions mirror biochemical profiles observed in classical dopamine pathway disorders and reinforce the plausibility of this targeted mechanism.

Importantly, individual genetic variation in viral entry receptors such as ACE2 and TMPRSS2, as well as dopamine-associated genes like DRD2 and DAT1, may modulate the vulnerability of dopaminergic neurons to SARS-CoV-2 infection. These genetic polymorphisms alter the surface expression of key receptors and transporters, affecting the susceptibility of neuronal populations to viral invasion. This variability likely underlies the broad spectrum of neurological manifestations observed among long COVID patients.

Taken together, the neuropathological, clinical, molecular, and genetic data converge to support a coherent model: SARS-CoV-2 may selectively target dopaminergic neurons in individuals with predisposing genetic expression profiles, producing region-specific neurotransmitter disruption that manifests as the heterogeneous symptomatology of long COVID.


\subsection*{Clinical Evidence of Dopaminergic Disruption}
Post-COVID patients show symptoms resembling dopamine depletion disorders \citep{brundin2020parkinsonmanagement, devine2022parkinsonscovid}. These include reduced dopamine transporter (DAT) availability, which correlates with fatigue and cognitive impairment; Parkinsonian features in some Long COVID patients, including tremor and bradykinesia; significant anhedonia and motivation deficits consistent with mesolimbic dopamine pathway dysfunction; and executive function impairments similar to those seen in ADHD and other dopamine-related disorders. Together, these observations suggest substantial dopaminergic involvement in Long COVID pathophysiology.

\subsection*{Genetic Variation in Relevant Surface Proteins}
Studies have identified high interindividual variability in ACE2, TMPRSS2, and DRD2 gene expression \citep{cao2020ace2genetics, hou2021ace2variants, eisenstein2016drd2}. This genetic heterogeneity could explain why only a subset of COVID-19 patients develop Long COVID symptoms, and why symptom profiles vary widely among those affected. Specific polymorphisms associated with altered dopamine signaling (such as DRD2 Taq1A and DAT1 VNTR variants) have been linked to motivational and cognitive traits that overlap with Long COVID symptoms \citep{gizer2009adhdgenes}.

\subsection*{Related Mechanisms and Pathways}
Beyond direct viral infection, several mechanisms may contribute to dopaminergic dysfunction in Long COVID. Mitochondrial impairment particularly affects high-energy neurons such as dopaminergic cells \citep{stefano2021mitochondrialfog}. Neuroinflammation can alter dopamine synthesis and signaling \citep{crunfli2023astrocytes}. Interactions with cholinergic systems, which are linked to attention and cognitive processing, may compound cognitive symptoms \citep{leitzke2025cholinergic}. Additionally, astrocytic dysfunction can impair neurotransmitter recycling and neuronal support \citep{crunfli2023astrocytes}. These mechanisms likely work in concert with direct viral effects to produce the complex symptomatology observed in Long COVID.

\begin{center}
\section*{Testable Predictions}
\end{center}

\subsection*{Genetic Predictions}
Long COVID patients are expected to show higher frequencies of specific ACE2, TMPRSS2, DRD2, and DAT1 polymorphisms compared to individuals who recovered without persistent symptoms. Genetic variants associated with increased ACE2 expression in dopaminergic regions should correlate with Long COVID severity. Furthermore, polygenic risk scores combining relevant dopaminergic and viral entry genes will likely predict Long COVID susceptibility better than individual variants. These genetic markers could potentially serve as diagnostic tools and guide personalized treatment approaches.

\subsection*{Neuroimaging Predictions}
PET imaging studies should reveal reduced dopamine transporter (DAT) availability in Long COVID patients compared to controls. Functional MRI would be expected to demonstrate altered activation patterns in dopamine-rich brain regions during reward and executive function tasks. The degree of dopaminergic abnormality on imaging should correlate with symptom severity, providing both a diagnostic marker and a potential means of monitoring treatment response. Such imaging findings would provide strong support for the dopaminergic hypothesis.

\subsection*{Experimental Predictions}
Cell culture experiments should show differential infection rates when exposing dopamine neurons with different genotypes to SARS-CoV-2. Animal models with humanized ACE2 expression would be expected to develop dopamine-related behavioral deficits following SARS-CoV-2 infection. Post-mortem brain tissue from Long COVID patients should show evidence of selective damage to dopaminergic neurons. These experimental approaches would provide direct evidence for the proposed mechanism of selective dopaminergic vulnerability in genetically susceptible individuals.

\begin{center}
\section*{Therapeutic Implications}
\end{center}

\subsection*{Restoring Dopamine Signaling}
Therapeutic strategies should focus on dopamine agonists such as pramipexole and ropinirole to directly stimulate dopamine receptors. DAT inhibitors including bupropion and methylphenidate could increase synaptic dopamine availability. L-DOPA therapy might enhance dopamine synthesis \citep{beauchamp2022parkinsonismwave, delprete2021parkinsoncovidimpact}. Neuroprotective agents could prevent further dopaminergic neuron loss \citep{limanaqi2022dopamineregulation, nataf2020ddcace2}. The selection of specific dopaminergic interventions would ideally be guided by symptom presentation and neuroimaging findings to target the most affected pathways.

\subsection*{Personalized Treatment Based on Genetic Profile}
Genetic screening could identify patients most likely to benefit from dopaminergic interventions. Treatment protocols could be tailored based on specific dopaminergic pathway involvement, as determined by symptom profiles and neuroimaging. Combined approaches addressing both dopaminergic dysfunction and neuroinflammation may be necessary for optimal symptom management. This personalized medicine approach would maximize therapeutic efficacy while minimizing unnecessary medication exposure in patients unlikely to benefit from dopaminergic interventions.

\begin{center}
\section*{Limitations and Alternative Explanations}
\end{center}
Several limitations and alternative explanations must be considered when evaluating this hypothesis. Long COVID likely involves multiple pathophysiological mechanisms beyond dopaminergic dysfunction, and the relative contribution of each remains to be determined. Neuroinflammation may cause dopaminergic symptoms without direct viral infection of neurons, acting instead through inflammatory mediators that disrupt dopamine signaling. The observed dopaminergic effects could be secondary to damage in other neural systems rather than representing the primary pathology. Additionally, the similarity to dopamine-related disorders could be coincidental rather than mechanistic. Further research is needed to determine the relative contribution of dopaminergic dysfunction compared to other proposed mechanisms and to establish whether it represents a primary or secondary phenomenon in Long COVID pathophysiology.

\begin{center}
\section*{Conclusion}
\end{center}
This hypothesis suggests that Long COVID's neurological effects may stem from direct viral infection of dopamine-producing neurons, but primarily in genetically susceptible individuals. The proposed mechanism—selective vulnerability of dopaminergic neurons due to genetic variability in surface receptor expression—offers an integrated explanation for the diverse neurological and cognitive symptoms observed in Long COVID. Importantly, this hypothesis generates specific testable predictions and suggests novel therapeutic approaches centered on dopamine restoration and neuroprotection. If supported by further research, this framework could lead to more effective, personalized treatments for Long COVID patients suffering from neurological symptoms.

\newpage
\bibliographystyle{plainnat}
\bibliography{references_complete_merged_final}

\end{document}
